\section{No Reading Reads Itself}

Place a specimen on an electronic balance. The display changes. That change is first a state of the device: these segments illuminate; these bits are stored. To say that the specimen has a certain mass requires more. The instrument must be responding to the quantity intended. Its scale and unit must be known. Its relation to a standard must have been established. The conditions must fall within its range, and the uncertainty must be judged. A displayed number may be an \emph{indication}; a measurement result is a warranted claim about the object.\endnote{See JCGM 200:2012, 2.9, 2.39, and 4.1; Stanford Encyclopedia of Philosophy, ``Measurement in Science,'' sec.~7.1; and NIST's Engineering Statistics Handbook, sec.~2.3.}

This does not make measurement arbitrary. It explains why measurement is rational. The result is more than the event in the machine because the event has been received into an order of meaning: a defined quantity, a model of the interaction, a chain of comparison, and an inference open to correction. The machine does indispensable work. It does not perform the whole act of knowing.

Even the purest-looking datum arrives with a question attached: \emph{a datum of what?} A pulse on a detector may be noise, malfunction, background, or the predicted trace of an event. More observations may decide among them, but only because each observation is interpreted under reasons for expecting one pattern rather than another. An absence of signal can be powerful evidence when a proposal predicts a signal under specified conditions. Apart from such a proposal, silence from this detector is only silence from this detector.

Here a pressure made famous by C.~S. Lewis can be used without asking it to bear more than it can. The physical cause of a belief and the rational ground for holding it are not interchangeable descriptions. A brain event may occur because another event caused it; a conclusion is known because the premises warrant it. Likewise, an instrument's state is caused by an interaction, while its standing as evidence depends upon reasons. The two accounts need not compete, but neither can replace the other.\endnote{The C.~S. Lewis Foundation's study guide identifies ch.~3 of \emph{Miracles} as Lewis's argument that naturalism is self-refuting and records his modification after Anscombe's criticism; the Foundation's chronology dates the revised reprint to 1960. Anscombe's published retrospective independently records the rewrite. The present cause--ground formulation and its application to measurement are project synthesis from those identified records: no Lewis wording is reproduced, the licensed revised chapter was not directly collated, and the analogy neither treats his broader argument as settled nor infers an immaterial mind.}

The point is almost embarrassingly close to the laboratory. Calibration is not another flash on the same unexamined display. It is a reasoned comparison. Error is not whatever the instrument dislikes. It is a judgment that the indication fails to correspond as it should to the quantity under investigation. The very practices that make instruments trustworthy therefore presuppose truth, reference, identity, valid inference, and the difference between seeming and being. None of these is an enemy of measurement. They are what allow a reading to be read.
